\documentclass[11pt]{article}

\usepackage[T1]{fontenc}
\usepackage[utf8]{inputenc}
\usepackage{amsmath}
\usepackage{amssymb}
\usepackage{amsthm}
\usepackage{mathtools}
\usepackage[hidelinks]{hyperref}
\usepackage{doi}

\setlength{\emergencystretch}{2em}

\newtheorem{proposition}{Proposition}
\newtheorem{definition}{Definition}
\newtheorem{hypothesis}{Hypothesis}
\newtheorem{remark}{Remark}

\newcommand{\Heis}{\mathrm{Heis}}
\newcommand{\Cay}{\mathrm{Cay}}
\newcommand{\Sent}{S_{\mathrm{ent}}}
\newcommand{\rpair}{r_{\mathrm{pair}}}
\newcommand{\spair}{\sigma^{\mathrm{can}}_{\mathrm{pair}}}
\newcommand{\Rset}{\mathcal{R}}

\title{Residual-Rank Entanglement of Conjugate Weil Pairs\\
in the Spectral Admissibility Cascade}
\author{J\'er\^ome Beau}
\date{June 2026}

\begin{document}

  \maketitle

  \begin{abstract}
    We compute the entanglement entropy of a closed conjugate Weil pair $\{c,q-c\}$ along the spectral
    admissibility cascade of the Cosmochrony programme, where entanglement is the infra-projectable
    invariant of a non-decomposable joint fibre of the non-injective projection.
    The bipartite state lives on the residual fibre-level support of the Gram--Schmidt span of Weil
    fingerprints on the Heisenberg Cayley graph.
    We show that the normalized admissibility residual of each newly admitted independent direction equals
    unity: the admitted Weil fingerprints are mutually orthogonal (proved analytically in the companion
    paper~\cite{Beau2026ent1}), so the reduced Schmidt spectrum is flat and the reduced entropy is exactly
    $\Sent(n)=\log \rpair(n)$, the logarithm of the residual Gram--Schmidt rank.
    No Born--Infeld spectral correction survives: the admissibility-weight deficit is identically zero.
    Because the cascade only adjoins independent directions, $\rpair(n)$ is non-increasing structurally,
    so the residual conjugate-pair entanglement is maximal at the earliest admissible stage and decreases
    monotonically as correlated capacity is stabilized into projected structure.
    The capacity-level statement (non-negativity of the stabilization rate) is established numerically.
    The entropy computed here is the full residual-rank observable on the complete Gram--Schmidt basis;
    its relation to the minimal admissible pair sector, whose entropy is $\log 2$ conditional on two
    explicit hypotheses (H-rank) and (H-inv), is established separately~\cite{Beau2026ent2}.
    The remaining finite-$q$ task is to measure the contraction law of $\rpair(n)$ and compare its exponent
    with the pair-capacity exponent $\delta_{\mathrm{pair}}$.
  \end{abstract}

  \noindent\textbf{Keywords:}
  non-injective projection;
  Weil representation;
  Heisenberg group;
  entanglement entropy;
  spectral admissibility;
  conjugate pairs;
  Schmidt rank;
  Born--Infeld parity.

  \newpage
  \tableofcontents

  \newpage
  \section{Introduction}
    \label{sec:intro}

    The Cosmochrony programme reconstructs physical structure from admissibility constraints on non-injective
    projections on the Heisenberg group $\Heis_3(\mathbb{Z}/q\mathbb{Z})$, with the non-injective projection
    $\Pi$ as the primitive of physical description~\cite{Cosmochrony,Beau2026m,Beau2026l}.
    Within this programme, entanglement is not an added postulate.
    It is the infra-projectable invariant of a non-decomposable joint fibre of $\Pi$: two degrees of freedom
    that share a single resolution cell of $\Pi$ cannot be assigned separate projected descriptions, and the
    correlation between them is exactly the part of the joint fibre that $\Pi$ cannot split.
    The same non-injective mechanism already yields the violation of Bell factorisation and the Tsirelson
    bound from the projection-fibre geometry~\cite{Beau2026b}, and the phase-coherent correlations of the
    quantum sector from the failure of factorisation under admissibility~\cite{Beau2026q1,Beau2026HeisenbergStructure}.

    What has been missing is a quantitative location of entanglement \emph{within} the spectral admissibility
    cascade itself: given the conjugate Weil structure that organises the cascade, how much entanglement does a
    conjugate pair carry, and how does it evolve as the cascade resolves?
    This note answers that question for a single closed conjugate pair $\{c,q-c\}$.

    The result is structurally simple.
    The reduced state of the pair, restricted to the residual fibre-level support of the Gram--Schmidt span of
    Weil fingerprints, has a flat Schmidt spectrum, because the admitted fingerprint directions are mutually
    orthogonal.
    Hence the reduced entropy is exactly the logarithm of the residual rank,
    \[
      \Sent(n)=\log \rpair(n),
    \]
    with no surviving spectral weighting.
    This entropy is maximal at the earliest admissible stage and decreases monotonically; the monotonicity is
    structural, following from the Gram--Schmidt construction alone.
    The mass hierarchy of the projected description is then the downstream trace of this stabilization, not an
    independent input~\cite{Beau2026a4,Beau2026c}.

    Two subsequent companion papers sharpen this result.
    The orthogonality of the admitted fingerprints, stated here as the property $w_j=1$, is proved
    analytically in~\cite{Beau2026ent1}: the $k{=}3$ Weil fingerprint is a discrete Fourier mode at frequency
    $A=c_1b_1+c_2b_2+c_3b_3 \bmod q$, and distinct modes are orthogonal.
    The relation between the present full residual-rank entropy and the minimal admissible pair sector is
    clarified in~\cite{Beau2026ent2}: under two explicit, independent hypotheses -- (H-rank), on the
    absence of additional effective symmetries beyond Born--Infeld parity, and (H-inv), on diagonal
    invariance of the proto-state -- the rank-3 admissible projection carries the spin-$\tfrac12$ singlet
    with entropy $\log 2$, while $\Sent(n)=\log\rpair(n)$ is the finer pipeline-level observable on the
    full Gram--Schmidt basis, unconditional for matched blocks.
    The two coincide only when $\rpair(n)=2$.

  \section{Conjugate Weil pairs and the bipartite fibre}
    \label{sec:pair}

    Let $\rho_c$ denote the Weil representation of $\Heis_3(\mathbb{Z}/q\mathbb{Z})$ attached to the central
    character $c$, acting on $V_c\cong\mathbb{C}^q$.
    For a nonzero central character $c$, the conjugate character is $q-c$, and the two representations are
    related by complex conjugation,
    \[
      V_{q-c}\simeq\overline{V_c},
    \]
    so $\{c,q-c\}$ forms a natural pair closed under the Born--Infeld parity involution.
    Born--Infeld parity forces the admissible fibre of the pair to be equatorial, fixing the eigenvalue ratio
    of the projected covariance and binding the two characters into a single admissible
    object~\cite{Beau2026a34,Beau2026a22}.

    The bipartite space at BFS depth $n$ is
    \[
      K_{c}(n)=W_c(n)\otimes W_{q-c}(n),
    \]
    where $W_c(n)$ is the admissible fibre-level subspace generated by $\rho_c$ up to depth $n$ (made precise in
    Section~\ref{sec:rank}).
    We treat the pair as a single closed object.

    \begin{hypothesis}[Closure of the conjugate pair]
      \label{hyp:closure}
      The conjugate pair $\{c,q-c\}$ is in a pure joint state on $K_c(n)$, with Schmidt vectors paired
      diagonally as $e_j\leftrightarrow\overline{e_j}$.
    \end{hypothesis}

    Hypothesis~\ref{hyp:closure} bundles two separate claims, and only one of them is derived.
    The bipartite composition $K_c(n) = W_c(n) \otimes W_{q-c}(n)$ and the purity of the joint state on it
    are a supplied composition structure (a K4 crossing~\cite{Beau2026esub}), not a consequence of
    admissibility or of the parity involution.
    What the parity involution $V_{q-c}\cong\overline{V_c}$ does derive, conditional on that supplied
    purity, is the pairing pattern: among purity-compatible Schmidt pairings, the diagonal one
    $e_j\leftrightarrow\overline{e_j}$ is the one selected, not an independent choice among possible
    pairings~\cite{Beau2026a34,Beau2026a21}.
    Hypothesis~\ref{hyp:closure} is the precise sense in which a conjugate pair is treated as a closed,
    infra-projectable object rather than two separately projected fibres; within that supplied
    composition, the reduced state of one factor is then well defined, and its entropy is the
    entanglement of the pair.

    \begin{remark}[Scope of E1's closure result]
      \label{rem:e1-scope}
      The companion paper E1~\cite{Beau2026ent1} proves the admission-orthogonality identity $w_j=1$ for
      admitted Weil fingerprint directions, upgrading the numerical hypothesis of this note
      (Section~\ref{sec:rank}) to a theorem, and additionally claims Hypothesis~\ref{hyp:closure} itself
      as a theorem, but only for matched single-character blocks $(c,0,0)$ and $(q-c,0,0)$: its argument
      identifies the residual subspaces $V^{(q-c)}_{\mathrm{res}}(n)=\overline{V^{(c)}_{\mathrm{res}}(n)}$
      from the Fourier-mode structure, then exhibits the maximally entangled state on their tensor
      product as ``the natural'' choice, rather than deriving purity from a more primitive input.
      By E1's own scope remark, this matched-block argument does not extend to the independently
      sampled blocks of the canonical pair observable used throughout this note: general or mismatched
      blocks $(c,c_2,c_3)$ and $(q-c,c_2',c_3')$, where either factor has $(c_2,c_3)\neq(0,0)$ or the two
      factors' auxiliary parameters differ.
      The Schmidt-partner pairing is not fixed in that case, and E1 leaves it open.
      Hypothesis~\ref{hyp:closure}, as used here for the general conjugate pair, therefore remains a
      supplied composition structure (a K4 crossing) independently of E1's restricted result.
    \end{remark}

  \section{The residual fibre-level rank}
    \label{sec:rank}

    The cascade observable is the projective capacity of a Weil block.
    For a generic block with central character $c$, breadth-first search on the Cayley graph
    $\Cay(\Heis_3(\mathbb{Z}/q\mathbb{Z}),S_q)$ generates, shell by shell, the Weil fingerprint vectors
    $\rho_c(g)v_0$, and a Gram--Schmidt pass records how many of them are new~\cite{Beau2026a15,Beau2026a16}.
    Writing $\delta r_n$ for the number of fingerprint vectors in shell $S_n$ that are linearly independent of
    the span built over shells $0,\dots,n-1$, the per-shell capacity is
    \[
      \sigma_c(n)=\frac{\delta r_n}{|S_n|}.
    \]
    The conjugate-pair observable is the product
    $\spair(n)=\sigma_c(n)\,\sigma_{q-c}(n)$, built at the level of the raw cumulative spans so that the
    residual amplitude factor $r(c,q)$ is treated as a normalization artifact rather than physical
    content~\cite{Beau2026a21,Beau2026a23}.

    \begin{definition}[Residual fibre-level rank]
      \label{def:rank}
      Let $R(n)=\sum_{m\le n}\delta r_m$ be the cumulative Gram--Schmidt span dimension at depth $n$, and let
      $R_\infty$ be its saturated value.
      The residual fibre-level rank is
      \[
        \rpair(n)=R_\infty-R(n),
      \]
      the number of independent admissible directions not yet activated at depth $n$.
    \end{definition}

    At complete locking $R(n)=R_\infty$, so $\rpair(n)=0$: this is the trivial separable endpoint.
    Throughout we read the entanglement entropy with the convention
    \[
      \Sent(n)=
      \begin{cases}
        \log \rpair(n), & \rpair(n)\ge 1,\\
        0,              & \text{at complete locking,}
      \end{cases}
    \]
    equivalently using the effective rank $\rpair^{\mathrm{eff}}(n)=\max(\rpair(n),1)$, so that the separable
    endpoint carries zero entanglement.

    Equivalently, tagging each admitted direction $j$ by its admission depth $m_j$, the residual support is
    $\Rset_n=\{\,j:m_j>n\,\}$ and $\rpair(n)=|\Rset_n|$.
    By Hypothesis~\ref{hyp:closure}, $\rpair(n)$ is the Schmidt rank of the conjugate pair at depth $n$ (the
    pair rank is the smaller of the two conjugate factors, which coincide by parity).

  \section{Reduced entropy: saturation of the rank bound}
    \label{sec:entropy}

    The admissibility-weight ansatz becomes trivial on the actual Gram--Schmidt data.
    For each newly admitted Weil fingerprint direction, the normalized residual at the moment of admission
    satisfies
    \[
      w_j=\frac{\|w_j^{\mathrm{GS}}\|}{\|v_j\|}=1
    \]
    in the tested finite-$q$ runs.
    Every admitted direction is thus exactly orthogonal to the previously generated span.
    This is not imposed by Gram--Schmidt, which would also admit a direction of partial overlap;
    it is a measured property of the Weil fingerprints, whose distinct images are mutually orthogonal until
    they become linearly dependent.
    The admissibility process is therefore binary at the level of independent directions:
    a fingerprint is orthogonal to the current span and admitted with $w_j=1$, or already contained in it and
    rejected.

    \begin{proposition}[Flat residual Schmidt spectrum]
      \label{prop:flat}
      If the admitted Weil fingerprint directions are mutually orthogonal, the normalized Schmidt weights on the
      residual support are flat,
      \[
        p_j(n)=\frac{1}{\rpair(n)},\qquad j\in\Rset_n,
      \]
      and the reduced entropy of the closed conjugate pair is exactly
      \[
        \Sent(n)
        =-\sum_{j\in\Rset_n}\frac{1}{\rpair(n)}\log\frac{1}{\rpair(n)}
        =\log \rpair(n).
      \]
      Equivalently the admissibility-weight deficit vanishes,
      \[
        \varepsilon_{\mathrm{adm}}(n):=\log \rpair(n)-\Sent(n)=0.
      \]
    \end{proposition}

    This equality is not an additional maximal-entanglement postulate.
    It follows from the orthogonality of the admitted Weil fingerprint directions in the Gram--Schmidt
    cascade.
    The identity $w_j=1$ was first verified across $q\in\{29,61,101,151\}$ and multiple generic blocks, and is
    now proved analytically in the companion paper~\cite{Beau2026ent1}: the $k{=}3$ Weil fingerprint is a
    discrete Fourier mode at frequency $A=c_1b_1+c_2b_2+c_3b_3 \bmod q$, and a newly admitted frequency is
    orthogonal to the entire prior span, so $w_j=1$ exactly.

  \section{Monotonicity of the entanglement profile}
    \label{sec:monotonicity}

    Since the Gram--Schmidt cascade only adjoins independent directions, $R(n)$ is non-decreasing.
    Therefore $\rpair(n)=R_\infty-R(n)$ is non-increasing.

    \begin{proposition}[Structural monotonicity]
      \label{prop:mono}
      The entanglement entropy $\Sent(n)=\log \rpair(n)$ is monotone non-increasing in $n$.
      The residual conjugate-pair entanglement is maximal at the earliest admissible stage and decreases as
      independent fingerprint directions are stabilized into the projected span.
    \end{proposition}

    This rank monotonicity is structural: it follows from the Gram--Schmidt construction itself, with no
    numerical input.
    A separate statement concerns the capacity-level stabilization rate
    \[
      \Delta I(n)=\spair(n)-\spair(n+1),
    \]
    which measures the rate at which residual correlated capacity is converted into projected stable structure.
    The non-negativity of $\Delta I(n)$ --- one-way activation at the level of the capacity --- is established
    numerically in the projective temporal ordering analysis, with no physically meaningful regressions
    observed; a general analytic proof remains open~\cite{Beau2026pto}.
    The monotonicity of the entanglement profile does not rely on this capacity-level result: it is already
    guaranteed by Proposition~\ref{prop:mono}.

  \section{Physical interpretation}
    \label{sec:interpretation}

    The conjugate-pair entanglement is maximal at the earliest admissible stage, where the residual rank is
    largest, and it decreases monotonically as the cascade resolves.
    There is no intermediate maximum: the trajectory is rank-dominated and, under the orthogonality of the
    admitted fingerprints, rank-determined.
    The decrease is not a loss of correlation into an environment; it is the conversion of residual correlated
    capacity into projected stable structure.
    The projected mass hierarchy is the downstream trace of this stabilization, recorded by the capacity
    exponent rather than by an independent dynamical input~\cite{Beau2026a4,Beau2026a25,Beau2026c}.

    The interpretation is consistent with the relational reading of time in the programme.
    A fully infra-projectable invariant has no position in the projected temporal ordering, which is itself the
    cumulative projected capacity~\cite{Beau2026pto}: it is a single unsplit object below the resolution of
    $\Pi$, not a pair of separately ordered events.
    Correlations of this kind therefore carry no signal and admit no influence between separated projected
    descriptions; locality at the projected level is preserved, in the same way that the non-injective
    mechanism reproduces the Tsirelson bound without superluminal signalling~\cite{Beau2026b}.
    Entanglement is, in this sense, the not-yet-time-ordered sector of the cascade.

  \section{Numerical deliverable}
    \label{sec:deliverable}

    The finite-$q$ calculation reduces to the residual fibre-level rank trajectory:
    \[
      \rpair(n),
      \qquad
      \Sent(n)=\log \rpair(n),
      \qquad
      \Delta I(n)=\spair(n)-\spair(n+1).
    \]
    For $q\in\{61,151,211,307\}$, the rank trajectory is reconstructed from the cumulative capacity profiles
    already produced by the systematic pair campaign~\cite{Beau2026a29}, via
    $\rpair(n)=R_\infty-\sum_{m\le n}\sigma_c(m)\,|S_m|$, and $\Delta I(n)$ follows directly from the stored
    $\spair$ profile.

    The task is therefore not to rescue an entropy correction --- there is none --- but to measure the
    contraction law of $\rpair(n)$, its scaling exponent, and its relation to the pair-capacity exponent
    $\delta_{\mathrm{pair}}$.
    If the contraction exponent agrees with $\delta_{\mathrm{pair}}$ within the expected finite-$q$ window, the
    entanglement trajectory and the capacity law share the same spectral engine.
    If it differs, the difference marks a genuine separation between residual rank contraction and shell-wise
    capacity decay.

  \section{Remark on possible non-rank corrections}
    \label{sec:remark}

    The equality $\Sent(n)=\log \rpair(n)$ uses the native Gram--Schmidt admissibility weights on Weil
    fingerprint directions, which are uniform.
    A non-trivial entropy deficit would require an additional weighting structure on the residual support, for
    example from the projected $\mathfrak{su}(2)$ coordinates of the admissible fibre or from an explicitly
    diagonalized Laplacian sector~\cite{Beau2026a31,Beau2026a33}.
    Such a correction would be a distinct modelling choice, defined on a different observable, and is not part
    of the present fibre-level entanglement calculation.
    That different observable is the minimal admissible pair sector of~\cite{Beau2026ent2}: the rank-3
    admissible projection, whose two-dimensional parity pair carries the spin-$\tfrac12$ singlet with
    entropy $\log 2$ under two explicit hypotheses, (H-rank) (spectral atomicity's general analytical
    proof) and (H-inv) (the singlet identification), neither derived in the papers that state them.
    It is distinct from the full residual-rank observable computed here.

  \section{Conclusion}
    \label{sec:conclusion}

    The entanglement of a closed conjugate Weil pair, computed honestly on the spectral admissibility cascade,
    returns to a pure O-series observable: the logarithm of the residual fibre-level rank.
    The chain is
    \[
      \begin{aligned}
        &\text{closed conjugate pair}+\text{orthogonal admitted fingerprints}\\
        &\qquad\Longrightarrow\ \text{flat Schmidt spectrum}\\
        &\qquad\Longrightarrow\ \Sent(n)=\log \rpair(n),
      \end{aligned}
    \]
    with monotone decrease guaranteed structurally by the Gram--Schmidt construction.
    No free parameter and no external weight enters.

    The epistemic status is stratified as follows.
    The reduction $\Sent(n)=\log \rpair(n)$ is derived, conditional on the orthogonality of the admitted
    fingerprints, which is numerically established for $q\in\{29,61,101,151\}$ and proved analytically in the
    companion paper~\cite{Beau2026ent1}.
    The monotonicity of $\Sent$ is structural.
    The capacity-level one-way activation governing $\Delta I$ is numerically established in the projective
    temporal ordering analysis, with a general analytic proof open.
    The purity of the conjugate pair (Hypothesis~\ref{hyp:closure}) is a supplied composition structure
    (a K4 crossing, Remark~\ref{rem:e1-scope} and Section~\ref{sec:pair}), not a consequence of the
    Born--Infeld parity involution; the parity involution derives only the diagonal pairing pattern
    conditional on that supplied purity, and the whole hypothesis remains the standing structural
    assumption of the construction.
    The entropy computed here is the full residual-rank observable; its separation from the minimal admissible
    pair sector ($\Sent^{\mathrm{adm}}=\log 2$ under the explicit hypotheses (H-rank) and (H-inv)) is treated
    in~\cite{Beau2026ent2}.

    \newpage
    \bibliographystyle{plain}
    \bibliography{cosmochrony-bibliography}

\end{document}
